\chapter*{Conclusion and Perspectives}
\addcontentsline{toc}{chapter}{Conclusion and Perspectives}

\section*{Conclusion}

This report presented the design, implementation, and evaluation of the
\textit{Knowledge Graph Exploration and Querying Platform with RAG}, developed
during an end-of-study internship at SCIENTIAE. The work contributed to the
development of a sovereign and domain-adaptable Graph-RAG framework through a
political and geopolitical application. Its primary objective was to transform
heterogeneous public information into a structured, searchable, traceable, and
explainable knowledge layer. The resulting platform was designed to support
factual, relational, temporal, and comparative analysis while preserving
explicit links between generated answers and their supporting evidence.

The project addressed a central limitation of conventional retrieval systems.
Although semantic search can identify passages related to a question, it does
not explicitly represent the relationships that connect actors, institutions,
events, and locations across several documents. Classical RAG improves answer
generation by providing external context to a language model, but its
passage-oriented retrieval remains insufficient for questions that require
entity disambiguation, temporal interpretation, or traversal across several
relations. The proposed approach therefore combines textual retrieval with
graph-based knowledge representation and reasoning.

The implemented platform covers the complete evidence-processing workflow.
Public information is acquired through RSS monitoring, query-driven
discovery, direct URL ingestion, and PDF parsing. Documents are subsequently
normalized and stored with their provenance before entity and relation
extraction. DSPy supports the optimization of this extraction stage, while
deterministic resolution rules reduce entity fragmentation. MongoDB preserves
source documents and processing artifacts, Neo4j stores canonical entities
and typed relations, and ZVec indexes textual chunks for semantic retrieval.
Together, these complementary representations provide the evidence required
for hybrid vector and graph retrieval.

The reasoning layer exposes three modes with different levels of analytical
depth. Standard RAG prioritizes direct retrieval and low response time.
Chain-of-thought reasoning decomposes complex questions and gathers evidence
over several steps. Reflexive reasoning additionally evaluates evidence
sufficiency and can initiate further retrieval before producing an answer.
The React analytical interface makes these capabilities available through
conversational querying, source consultation, graph exploration, specialized
political views, and administrative supervision. Furthermore, containerized
deployment and observability services provide operational visibility across
ingestion, retrieval, reasoning, and generation.

The system was evaluated using a curated benchmark of \textbf{150 questions}
covering factual, relational, temporal, and comparative tasks. Each question
was executed under the three reasoning modes, resulting in \textbf{450
executions}. The findings demonstrate a progressive improvement in answer
quality as reasoning depth increases. Average answer correctness rose from
\textbf{0.52} for standard RAG to \textbf{0.65} for chain-of-thought reasoning
and \textbf{0.80} for reflexive reasoning. The reflexive mode also achieved
the strongest overall results for answer relevancy, faithfulness, context
precision, and context recall. However, these improvements were accompanied
by higher latency because additional retrieval and verification operations
were required.

The embedding comparison also demonstrated the importance of the retrieval
configuration. The Mistral embedding service reduced average response time
from \textbf{15 to 5 seconds} for standard RAG, from \textbf{30 to 15 seconds}
for chain-of-thought reasoning, and from \textbf{90 to 50 seconds} for
reflexive reasoning. Retrieval-quality indicators improved at the same time.
At the ingestion stage, DSPy optimization increased the composite extraction
score from \textbf{0.723} to \textbf{0.883} and reduced mean processing
latency from \textbf{4.62} to \textbf{3.21 seconds}. These results confirm
that Graph-RAG quality depends not only on answer generation but also on
corpus preparation, extraction reliability, knowledge representation, and
evidence retrieval.

Several limitations must nevertheless be considered when interpreting these
results. Evaluation was performed on a fixed corpus snapshot and a curated
benchmark; consequently, the reported scores cannot represent every
geopolitical region, language, event type, or analytical formulation.
Political information also evolves rapidly, which can affect the temporal
validity of extracted relations and reference answers. In addition,
LLM-assisted evaluation may introduce judge bias, while incomplete source
coverage, unsuitable chunk boundaries, ambiguous entities, and syndicated
content can limit retrieval independently of the selected reasoning mode.
These limitations do not invalidate the observed results, but they define the
conditions under which the conclusions should be understood.

Overall, the project demonstrates the feasibility of an inspectable Graph-RAG
framework that integrates acquisition, extraction, structuring, retrieval,
reasoning, and evaluation within a coherent platform. Its political and
geopolitical implementation shows that unstructured information can be
transformed into connected evidence and grounded answers with explicit source
attribution. More broadly, the work confirms that dependable Graph-RAG
behavior is a system-level property that depends on the combined quality of
data, models, storage structures, retrieval strategies, operational controls,
and evaluation procedures.

\section*{Perspectives}

Although the developed platform reached a functional and evaluable state,
several research and engineering directions remain open. The first direction
concerns retrieval quality. Future work should investigate graph-aware
reranking, stronger query expansion, temporal filtering, and more systematic
chunking strategies. These improvements could increase both context recall
and precision. Re-embedding the corpus with higher-dimensional multilingual
models would also support the extension of the platform to French and Arabic
political content.

A second direction concerns knowledge-graph enrichment. Future versions could
introduce more precise temporal validity intervals, improved event
normalization, confidence scores for extracted facts, and more advanced entity
disambiguation. Such improvements would increase graph reliability for
longitudinal political analysis. They would also support questions concerning
institutional evolution, alliances, conflicts, and changes in political
positions over time.

A third perspective concerns evaluation. Although the 150-question benchmark
supports comparison across question categories and reasoning modes, it cannot
represent every geopolitical region, language, event type, or analytical
formulation. Future evaluation should incorporate human analyst ratings,
additional multilingual and temporal questions, confidence calibration, and
adversarial cases requiring abstention or careful qualification.

Finally, continuous ingestion monitoring, richer administrative dashboards,
improved observability, and deployment hardening could strengthen the
platform operationally. These improvements would facilitate the transition
from a functional prototype to a robust decision-support tool. In particular,
they would improve its maintainability, auditability, and extensibility in a
production environment.
